点估计给出一个中心，却没有交代这个中心离真值可能有多远。区间估计把推断的输出从一个数扩成一段范围，写进报告只多几个字符，读错的方式却比点估计多得多。最常见的一种读法是把 \(95\%\) 当成``参数在这段里''的把握，而这句话在频率学派的框架里没有定义——它属于另一套框架，需要先给参数配一个分布才谈得上。

第一件事是把概率安在正确的对象上。频率学派的区间是一套程序：数据随机，端点随之随机，参数从头到尾不动；\(95\%\) 描述的是这套程序在重复使用下盖住真值的比例，而不是你手上这一条区间的属性。这个区分决定了报告该怎么措辞，也决定了拿到一个区间之后第一个该问的问题——不是端点有多宽，而是这套程序在你的样本量和参数范围内还有没有那个标称覆盖。Wald 区间在比例接近边界时的失灵，是这个问题的标准反面教材。这是\hyperref[sec:11-Mathematical-Statistics/04-Interval-Estimation/01]{``置信度属于程序，而非这次参数''}。

第二件事是把区间造出来。通用的做法是找一个形式上含参数、分布中不含参数的量，先在它的已知分布上截出一段概率，再把参数从不等式里反解出来。正态模型里这条路给出精确的 \(z\) 区间与 \(t\) 区间，也解释了方差区间为什么天然不对称。找不到这种精确的量时按顺序退让：退到只在大样本下成立的版本，再退到用重采样把分布模拟出来。这是\hyperref[sec:11-Mathematical-Statistics/04-Interval-Estimation/02]{``枢轴量把未知参数隔离出来''}。

第三件事是把两套语言并排看清。大样本下中心极限定理支撑起通用的正态近似区间，而 Bayesian 的可信区间直接对后验分布取概率，它的表述恰好就是第一节禁止频率区间说的那句话。正则条件下两者的端点常常近到看不出差别，但一个数的是重复实验里的比例，另一个量的是后验密度下的面积，接近不等于同义。这是\hyperref[sec:11-Mathematical-Statistics/04-Interval-Estimation/03]{``渐近区间与 Bayesian 可信区间''}。

有一条约定贯穿全单元：置信水平 \(1-\alpha\) 是程序的标称值，实际覆盖率是需要另行核实的事实，两者相等从来不是自动的。区间宽度则在收数据之前就基本定了，半宽按 \(1/\sqrt n\) 收缩，把它写进实验设计的成本，远低于事后发现区间太宽的代价。
